% ============================================================
% NBA Player Prop Prediction — Research Paper
% ============================================================
% Suggested venue: PLOS ONE or MIT Sloan Sports Analytics
% Compile: pdflatex paper.tex  (run twice for references)
% Bibliography: bibtex paper   (then pdflatex twice more)
% ============================================================

\documentclass[12pt, letterpaper]{article}

% ── Packages ─────────────────────────────────────────────────────────────────
\usepackage[margin=1in]{geometry}
\usepackage{amsmath, amssymb}
\usepackage{graphicx}
\usepackage{booktabs}           % professional tables (\toprule, \midrule, \bottomrule)
\usepackage{multirow}           % multi-row table cells
\usepackage{array}              % extended column formats
\usepackage{hyperref}           % clickable citations / URLs
\usepackage{url}
\usepackage{natbib}             % author-year citations: \citep{} \citet{}
\usepackage{caption}
\usepackage{subcaption}         % side-by-side figures
\usepackage{setspace}
\usepackage{parskip}            % space between paragraphs instead of indent
\usepackage{xcolor}
\usepackage{listings}           % code snippets if needed
\usepackage{float}              % [H] figure placement

\hypersetup{
    colorlinks=true,
    linkcolor=blue,
    citecolor=blue,
    urlcolor=blue
}

\doublespacing                  % remove for submission if venue requires single

% ── Title block ───────────────────────────────────────────────────────────────
\title{
    \textbf{Gradient Boosting Regression for NBA Player Prop Prediction:}\\
    \large Walk-Forward Validation, Feature Attribution,\\
    and Simulated Betting Performance
}

\author{
    Abba Otieno Ndomo\\
    \small [Institution]\\
    \small \href{mailto:abba.ndomo@gmail.com}{abba.ndomo@gmail.com}
}

\date{\today}

% =============================================================
\begin{document}
% =============================================================

\maketitle

% ── Abstract ──────────────────────────────────────────────────────────────────
\begin{abstract}
Player proposition (prop) betting markets in the NBA present a challenging
regression problem: predicting individual player statistics --- points, rebounds,
and assists --- from historical performance and contextual features.
We develop a gradient boosting regression system that trains separate
XGBoost models for each statistic using per-player rolling averages,
exponential moving averages, season baselines, hot/cold streak indicators,
and opponent defensive context features.
Trained on 1.1 million player-game observations spanning 2016--2024,
the models are evaluated under walk-forward cross-validation across four
held-out NBA seasons (2021--2024), yielding mean absolute errors of
4.84 points, 1.97 rebounds, and 1.41 assists.
We compare performance against three baselines --- a naive recency heuristic,
linear regression, and random forest --- all trained on identical feature sets.
SHAP (SHapley Additive exPlanations) analysis identifies the 10-game rolling
average, season expanding mean, and exponential moving average as the three
most influential predictors across all statistics, while opponent defensive
context contributes meaningful secondary signal.
A simulated betting experiment using synthetic market lines demonstrates
that the XGBoost model generates positive return on investment relative to
a linear regression baseline (+13.1\% ROI on points at a 1-point edge
threshold, $N=17{,}685$ bets), suggesting the model captures non-linear
structure beyond what simpler models extract from the same features.
We discuss limitations of synthetic line evaluation and directions for
future work using historical closing lines.

\medskip
\noindent\textbf{Keywords:} NBA, player proposition betting, gradient boosting,
XGBoost, SHAP, walk-forward cross-validation, sports analytics
\end{abstract}

\newpage
\tableofcontents
\newpage

% =============================================================
\section{Introduction}
\label{sec:intro}
% =============================================================

The NBA player proposition (prop) betting market has grown substantially
following the 2018 repeal of the Professional and Amateur Sports Protection
Act (PASPA), with per-game player prop volume now rivalling traditional
moneyline and spread markets~\citep{galekwa2024systematic}.
Sportsbooks set prop lines by estimating expected statistical output ---
typically points, rebounds, or assists --- for individual players, and bettors
wager on whether a player will exceed (\textit{over}) or fall short
(\textit{under}) of that line.
For a bettor, consistently predicting these values more accurately than the
market line --- even marginally --- can yield positive expected value at the
standard -110 juice, which requires only a 52.38\% win rate to break even.

The academic literature on basketball prediction has focused predominantly on
binary game-outcome classification (win/loss)~\citep{chen2021hybrid,ouyang2024xgboost,nature2025stacked}.
Individual player stat prediction, which requires regression over a continuous
target rather than a binary one, has received comparatively less
attention~\citep{galekwa2024systematic, plosone2025systematic}.
This distinction matters practically: game-outcome models are not directly
actionable for prop bettors, who need a point estimate and uncertainty
bound for a single player's output on a given night.
Furthermore, most prior work evaluates on a single chronological split,
which overstates generalization performance when the data distribution
shifts across seasons due to changes in pace, lineup construction,
and rule enforcement.

This paper makes the following contributions:
\begin{enumerate}
    \item A gradient boosting regression pipeline that trains separate
          XGBoost~\citep{chen2016xgboost} models for points, rebounds,
          and assists using stat-specific feature sets incorporating
          rolling windows, exponential smoothing, and opponent defensive context.
    \item A walk-forward cross-validation protocol that evaluates temporal
          generalization across four held-out NBA seasons (2021--2024),
          avoiding the look-ahead bias of a single chronological
          split~\citep{cerqueira2020evaluating}.
    \item Comparison against three baselines --- a naive last-5-game average,
          linear regression, and random forest --- all on identical features,
          isolating the contribution of model complexity from feature
          engineering.
    \item SHAP~\citep{lundberg2017unified} feature attribution identifying
          which features drive predictions and validating that opponent
          defensive context and exponential momentum add signal beyond
          simple recency.
    \item A simulated betting experiment demonstrating positive return on
          investment relative to a linear regression synthetic line,
          with explicit discussion of limitations inherent to any
          evaluation against synthetic rather than market-closing lines.
\end{enumerate}

% =============================================================
\section{Related Work}
\label{sec:related}
% =============================================================

\subsection{Basketball Outcome Prediction}

Machine learning has been applied extensively to NBA game-outcome
classification.
\citet{chen2021hybrid} proposed a two-stage XGBoost model for NBA game
outcome prediction, identifying lagged shooting percentages, defensive
rebounds, and assists as the most predictive features.
\citet{ouyang2024xgboost} extended this with SHAP analysis on 2021--2023
NBA data, finding field goal percentage and defensive rebounds to be
consistently dominant at the team level.
\citet{nature2025stacked} demonstrated that stacked ensemble approaches
combining multiple base learners improved game-outcome accuracy over
any single model.
In contrast to these game-level classifiers, our work targets individual
player statistics --- a continuous regression task with higher variance,
as player-level output is subject to foul trouble, minutes management,
and in-game matchup adjustments that cancel out at the team level.

\subsection{Sports Betting Market Efficiency}

\citet{baryla2007learning} demonstrated statistically significant
inefficiencies in NBA totals lines early in the season, achieving a
56.72\% win rate over 20 years by exploiting early-season bias in
closing totals, suggesting that market participants require several weeks
of current-season data to form accurate priors.
A systematic review by \citet{galekwa2024systematic} notes that
machine learning approaches to sports betting consistently outperform
naive baselines in backtesting but that the magnitude of edge shrinks
substantially when evaluated against sharp closing lines rather than
opening or mid-market prices.
Our betting simulation is explicitly positioned as a relative benchmark
(XGBoost vs.\ linear regression on identical features) rather than a
claim of absolute profitability against real sportsbook lines.

\subsection{Gradient Boosting for Sports Prediction}

XGBoost~\citep{chen2016xgboost} has become a standard benchmark on
tabular prediction tasks due to its handling of mixed feature types,
built-in regularization, and computational efficiency.
It has been applied successfully to game-level basketball
prediction~\citep{chen2021hybrid,ouyang2024xgboost} and outperforms
deep learning approaches on feature-engineered tabular data in
most sports analytics benchmarks~\citep{plosone2025systematic}.
Our use of per-player rolling features and opponent context is
consistent with the feature engineering approach reported as most
effective in this literature.

\subsection{Explainability in Sports Models}

SHAP~\citep{lundberg2017unified} provides a game-theoretically grounded
decomposition of model predictions into per-feature additive contributions,
enabling consistent comparison of feature importance across models and
data subsets.
\citet{ouyang2024xgboost} applied SHAP to a game-outcome XGBoost model
and found field goal percentage and team defensive rating to be the most
influential team-level features.
To our knowledge, SHAP has not previously been applied to individual
player stat regression for prop betting, making our feature attribution
analysis a novel contribution to this sub-field.

% =============================================================
\section{Data and Feature Engineering}
\label{sec:data}
% =============================================================

\subsection{Dataset}

The primary dataset comprises 1,648,656 player-game records spanning NBA
seasons from 1951 to 2025, sourced from the ESPN unofficial public API.
After filtering to Regular Season, Playoff, and Play-in Tournament games
with a minimum of 10 minutes played, and restricting to seasons from 2016
onward to reflect modern-era play style, 1,120,817 records remain across
4,273 unique players.
The 10-minute minimum excludes garbage-time appearances and players with
insufficient context to build meaningful rolling features.
The 2016 cutoff reflects the rapid structural shift in NBA offence
following the emergence of the three-point revolution; including earlier
data risks introducing distribution mismatch between training and test.

\subsection{Per-Player Rolling Features}

For each player, the following features are computed using a causal rolling
window (shifted by one game to prevent data leakage):

\begin{itemize}
    \item \textbf{Rolling means:} $\bar{x}^{(5)}$, $\bar{x}^{(10)}$ ---
          5- and 10-game simple moving averages for points, rebounds,
          assists, minutes, and field goal percentage.
    \item \textbf{Exponential moving average:}
          $\text{EMA}^{(5)}_t = \alpha x_{t-1} + (1-\alpha)\text{EMA}^{(5)}_{t-1}$,
          with span $s=5$ ($\alpha = 2/(s+1)$), capturing recent momentum
          with exponentially decaying weights.
    \item \textbf{Rolling standard deviation:} $\sigma^{(10)}$ over the
          last 10 games, used as a dispersion estimate for probability
          calculation.
    \item \textbf{Days rest:} days since the player's previous game,
          clipped at 10.
    \item \textbf{Season expanding average:} $\mu^{\text{season}}_t =
          \frac{1}{t-1}\sum_{k=1}^{t-1} x_k$, computed per player per
          season year.
    \item \textbf{Hot/cold streak:}
          \begin{equation}
              \text{hot\_cold} = \frac{\bar{x}^{(5)} - \mu^{\text{season}}}
              {|\mu^{\text{season}}| + 0.1}
              \label{eq:hotcold}
          \end{equation}
          A positive value indicates the player is performing above their
          season average over the last 5 games.
\end{itemize}

\subsection{Opponent Defensive Context}

For each game, player statistics are aggregated at the team level by
grouping on (game identifier, opposing team). This yields the total points,
rebounds, and assists allowed by each defending team per game.
A 10-game rolling average of these totals, shifted by one game
($\text{opp\_pts\_allowed\_L10}$, etc.), is then merged back to each
player row as a defensive context feature.
This approach approximates the ``defensive rating'' concept used in
basketball analytics but is computed from the same player-level data
as the target variable, avoiding the need for external box-score aggregates.

\subsection{Feature Sets by Statistic}

Each statistic uses a distinct feature set reflecting the different
contextual factors that drive scoring, rebounding, and playmaking.
Table~\ref{tab:features} lists the full feature sets.

\begin{table}[H]
\centering
\caption{Feature sets by prediction target.}
\label{tab:features}
\begin{tabular}{llp{7cm}}
\toprule
\textbf{Stat} & \textbf{Features} & \textbf{Notes} \\
\midrule
\multirow{3}{*}{Points}
  & is\_home, days\_rest & Situational \\
  & pts\_L5/10, pts\_ema\_L5, pts\_std\_L10 & Recency \\
  & reb\_L5, ast\_L5, min\_L5/10, fg\_pct\_L5/10 & Role context \\
  & season\_avg\_pts, hot\_cold\_pts & Baseline / form \\
  & opp\_pts\_allowed\_L10 & Opponent defense \\
\midrule
\multirow{3}{*}{Rebounds}
  & is\_home, days\_rest & Situational \\
  & reb\_L5/10, reb\_ema\_L5, reb\_std\_L10 & Recency \\
  & pts\_L5, min\_L5/10 & Role context \\
  & season\_avg\_reb, hot\_cold\_reb & Baseline / form \\
  & opp\_reb\_allowed\_L10 & Opponent defense \\
\midrule
\multirow{3}{*}{Assists}
  & is\_home, days\_rest & Situational \\
  & ast\_L5/10, ast\_ema\_L5, ast\_std\_L10 & Recency \\
  & pts\_L5, min\_L5/10, fg\_pct\_L5 & Role context \\
  & season\_avg\_ast, hot\_cold\_ast & Baseline / form \\
  & opp\_ast\_allowed\_L10 & Opponent defense \\
\bottomrule
\end{tabular}
\end{table}

% =============================================================
\section{Methodology}
\label{sec:method}
% =============================================================

\subsection{Model Architecture}

Three independent XGBoost regressors~\citep{chen2016xgboost} are trained,
one per statistic, using the squared error objective
($\texttt{reg:squarederror}$).
Hyperparameters are held constant across all three models:
learning rate $\eta = 0.05$, maximum tree depth $= 6$,
subsample ratio $= 0.8$, column subsample $= 0.8$,
minimum child weight $= 5$.
Early stopping with a patience of 50 rounds prevents overfitting,
with the test fold used as the validation set during training.
The same architecture is used in all four walk-forward folds and in
the final production model.

\subsection{Walk-Forward Cross-Validation}
\label{subsec:wfcv}

To evaluate temporal generalization, we employ a walk-forward
(rolling-origin) cross-validation scheme~\citep{cerqueira2020evaluating}.
Training data grows by one season per fold, and the model is tested
on the immediately following season only:

\begin{table}[H]
\centering
\caption{Walk-forward fold definitions.}
\label{tab:folds}
\begin{tabular}{ccc}
\toprule
\textbf{Fold} & \textbf{Train} & \textbf{Test} \\
\midrule
1 & 2016--2020 & 2021 \\
2 & 2016--2021 & 2022 \\
3 & 2016--2022 & 2023 \\
4 & 2016--2023 & 2024 \\
\bottomrule
\end{tabular}
\end{table}

This design mirrors real deployment: a model retrained at the start of each
season on all prior data tests on that season's games.
Test folds contain between 20,338 (2023) and 28,338 (2021) player-game
observations.

\subsection{Baseline Models}

Four models are compared under identical walk-forward folds:

\begin{enumerate}
    \item \textbf{Naive L5}: predicts the player's last-5-game rolling
          average. Requires no training and represents a pure recency
          heuristic --- effectively what a bettor would use without any
          model.
    \item \textbf{Linear Regression}: ordinary least squares on the full
          feature set for each statistic.
    \item \textbf{Random Forest}: an ensemble of 200 decision trees with
          maximum depth 10, minimum samples per leaf 5, trained with a
          fixed random seed for reproducibility.
    \item \textbf{XGBoost}: gradient boosted trees as described above.
\end{enumerate}

All models receive the same feature matrix; no additional features are
made available to any model. This isolates the effect of model capacity
from the effect of feature engineering.

\subsection{SHAP Feature Attribution}

SHAP (SHapley Additive exPlanations)~\citep{lundberg2017unified} values
are computed using \texttt{shap.TreeExplainer} on a random sample of
5,000 player-game observations from the 2024 season hold-out, using
the final model trained on all prior seasons (2016--2023).
Mean absolute SHAP values rank features by their average contribution
to prediction magnitude.
SHAP values are additive: the prediction for any observation equals the
model's base value plus the sum of all SHAP values, providing a
faithful attribution without the double-counting issues of Gini importance
or permutation methods.

\subsection{Betting Simulation}
\label{subsec:sim}

A flat-unit betting simulation is conducted over all walk-forward
test predictions (2021--2024, $N = 94{,}462$ player-game observations
per statistic).

\paragraph{Synthetic market lines.}
In the absence of historical sportsbook closing lines, we construct
three synthetic line proxies of increasing sophistication:
(i) the player's 10-game simple rolling average ($L_{10}$);
(ii) a 5-game exponential moving average ($\text{EMA}_5$);
and (iii) the linear regression model's prediction for that game,
using identical features to XGBoost.
The third proxy is the most conservative benchmark: both models see
the same information, so any positive XGBoost edge reflects
non-linear model capacity rather than information advantage.
The EMA$_5$ line over-weights the most recent 1--2 games and is
therefore \textit{noisier} than the $L_{10}$ line, not more efficient;
positive ROI against it does not constitute evidence of market
inefficiency.

\paragraph{Bet signal.}
A bet is placed when the absolute difference between the XGBoost
projection and the synthetic line exceeds a threshold $\tau$:
\begin{equation}
    \text{Bet OVER if } \hat{y}_{\text{XGB}} > \ell + \tau, \quad
    \text{Bet UNDER if } \hat{y}_{\text{XGB}} < \ell - \tau
    \label{eq:betsignal}
\end{equation}
where $\ell$ is the synthetic line value and $\tau \in \{0.0, 0.5,
1.0, 1.5, 2.0, 2.5\}$ (in stat units).

\paragraph{Payout structure.}
Standard -110 juice is assumed: a winning bet returns $+0.909$ units
per unit risked; a losing bet costs $-1.0$ unit.
The break-even win rate is $52.38\%$. ROI is defined as total profit
divided by number of bets placed.

% =============================================================
\section{Results}
\label{sec:results}
% =============================================================

\subsection{Prediction Accuracy}

Table~\ref{tab:mae} reports mean absolute error (MAE) averaged across
all four walk-forward folds for each model and statistic.
Full per-fold results appear in Appendix~\ref{app:folds}.

\begin{table}[H]
\centering
\caption{Walk-forward MAE by model and statistic (mean across folds 2021--2024).
         Bold indicates best performance per stat.}
\label{tab:mae}
\begin{tabular}{lccccc}
\toprule
\textbf{Stat} & \textbf{Naive L5} & \textbf{Linear} & \textbf{Rand.\ Forest} & \textbf{XGBoost} & \textbf{$\Delta$ vs Naive} \\
\midrule
Points    & 5.082 & 4.884 & 4.839 & \textbf{4.840} & $-4.8\%$ \\
Rebounds  & 2.070 & 1.986 & \textbf{1.968} & 1.966 & $-5.0\%$ \\
Assists   & 1.478 & 1.428 & 1.413 & \textbf{1.412} & $-4.5\%$ \\
\bottomrule
\end{tabular}
\end{table}

XGBoost improves over the naive L5 heuristic by 4.5--5.0\% in mean MAE
across all three statistics.
The improvement over linear regression is more modest (0.9--1.2\%),
which indicates that the bulk of the gain over the naive baseline is
attributable to the richer feature set rather than to model non-linearity.
Notably, random forest and XGBoost perform almost identically: their MAEs
differ by at most 0.002 points, 0.002 rebounds, and 0.001 assists.
This near-parity suggests the dataset is well-described by additive
non-linear interactions that both ensemble methods capture, and that
marginal gains from XGBoost's sequential boosting over random forest's
parallel bagging are negligible at this feature dimensionality.

Fold-level results (Appendix~\ref{app:folds}) show consistent performance
across years with a slight upward trend in MAE from fold 1 (2021) to fold 4
(2024) for all models.
For points, XGBoost MAE increases from 4.770 (2021) to 4.859 (2024),
an increase of 1.9\% over four seasons.
This drift is observed equally in all four models, suggesting it reflects
increasing variance in the test data --- potentially attributable to
roster turnover, pace changes, or load management --- rather than
model degradation.

\subsection{SHAP Feature Importance}

Figure~\ref{fig:shap} shows SHAP beeswarm plots for all three models on a
5,000-observation random sample from the 2024 test set.
Full ranked importance tables appear in Appendix~\ref{app:shap}.

\begin{figure}[H]
    \centering
    \begin{subfigure}[b]{0.32\textwidth}
        \includegraphics[width=\textwidth]{figures/shap_beeswarm_pts.png}
        \caption{Points}
    \end{subfigure}
    \hfill
    \begin{subfigure}[b]{0.32\textwidth}
        \includegraphics[width=\textwidth]{figures/shap_beeswarm_reb.png}
        \caption{Rebounds}
    \end{subfigure}
    \hfill
    \begin{subfigure}[b]{0.32\textwidth}
        \includegraphics[width=\textwidth]{figures/shap_beeswarm_ast.png}
        \caption{Assists}
    \end{subfigure}
    \caption{SHAP beeswarm plots for the final XGBoost models evaluated on
             5,000 randomly sampled 2024-season observations. Each dot
             represents one player-game; colour encodes feature value
             (red = high, blue = low); horizontal position encodes SHAP
             contribution to the prediction.}
    \label{fig:shap}
\end{figure}

\paragraph{Points model.}
The 10-game rolling average (\texttt{pts\_L10}) is the dominant feature
by a large margin (mean $|\text{SHAP}| = 2.752$), followed by the season
expanding mean (\texttt{season\_avg\_pts} = 1.134) and the 5-game
exponential moving average (\texttt{pts\_ema\_L5} = 1.100).
The near-equal contribution of \texttt{season\_avg\_pts} and
\texttt{pts\_ema\_L5} confirms that both the player's long-run baseline
and recent momentum carry independent predictive signal.
Opponent defensive context (\texttt{opp\_pts\_allowed\_L10} = 0.346) ranks
fourth, validating that matchup difficulty adds information beyond
the player's own history.
Minutes played (\texttt{min\_L5} = 0.313) ranks fifth, reflecting
that role fluctuations are a meaningful source of variance in scoring.
Home/away status (\texttt{is\_home} = 0.187) and days rest (0.145) add
modest but non-negligible contributions.

\paragraph{Rebounds model.}
The same three-feature hierarchy appears: \texttt{reb\_L10} (0.931),
\texttt{season\_avg\_reb} (0.643), \texttt{reb\_ema\_L5} (0.367).
Minutes features (\texttt{min\_L5} = 0.185, \texttt{min\_L10} = 0.159)
rank higher relative to the points model, consistent with the fact that
rebounding opportunity is more directly tied to time on the floor than
scoring, which is also influenced by shot selection and usage.
Opponent rebounds allowed (\texttt{opp\_reb\_allowed\_L10} = 0.075)
ranks seventh --- lower than for points --- suggesting defensive context
has less marginal value for rebounding prediction, possibly because
team rebounding rates are less variable than points allowed per game.

\paragraph{Assists model.}
The top-three structure is again \texttt{ast\_L10} (0.770),
\texttt{season\_avg\_ast} (0.406), \texttt{ast\_ema\_L5} (0.247).
A distinctive finding for assists is that \texttt{is\_home} ranks
fourth (0.093), ahead of opponent defensive context (0.089).
This is consistent with the well-documented home-court advantage in
playmaking: point guards and facilitators typically produce more assists
at home, where crowd noise may disrupt opposing defensive communication
and where travel fatigue does not impair decision-making.
Opponent assists allowed (\texttt{opp\_ast\_allowed\_L10} = 0.089)
ranks fifth, suggesting that teams that allow high assist totals to
opponents provide a favourable matchup for playmakers.

\paragraph{Cross-stat summary.}
In all three models, the top-3 features follow the same hierarchy:
$L_{10}$ rolling mean $>$ season expanding mean $>$ EMA$_5$.
The consistent third-place ranking of EMA$_5$ --- despite it being
derived from a shorter window than $L_{10}$ --- confirms that
exponentially decayed recent momentum carries information orthogonal
to the simple 10-game average, justifying its inclusion as a
non-redundant feature.

\subsection{Betting Simulation}

Table~\ref{tab:betting} reports simulated ROI at representative edge
thresholds across all three synthetic line definitions.
Full results across all thresholds appear in Appendix~\ref{app:betting}.

\begin{table}[H]
\centering
\caption{Simulated ROI by statistic, line type, and edge threshold
         ($N$ = number of bets; break-even ROI = $-0.048$ at $-110$ juice).
         Only thresholds with $N \geq 1{,}000$ reported.}
\label{tab:betting}
\begin{tabular}{llccc}
\toprule
\textbf{Stat} & \textbf{Line} & \textbf{Threshold} & \textbf{N bets} & \textbf{ROI} \\
\midrule
\multirow{3}{*}{Points}
  & $L_{10}$  & 1.0 & 42,918 & $+0.143$ \\
  & EMA$_5$   & 1.0 & 55,410 & $+0.198$ \\
  & LinReg    & 1.0 & 17,685 & $+0.131$ \\
\midrule
\multirow{3}{*}{Rebounds}
  & $L_{10}$  & 1.0 &  6,416 & $+0.250$ \\
  & EMA$_5$   & 1.0 & 17,559 & $+0.342$ \\
  & LinReg    & 1.0 &  1,525 & $+0.208$ \\
\midrule
\multirow{3}{*}{Assists}
  & $L_{10}$  & 1.0 &  2,552 & $+0.295$ \\
  & EMA$_5$   & 1.0 &  7,985 & $+0.339$ \\
  & LinReg    & 1.0 &    793 & $+0.156$ \\
\bottomrule
\end{tabular}
\end{table}

\begin{figure}[H]
    \centering
    \begin{subfigure}[b]{0.32\textwidth}
        \includegraphics[width=\textwidth]{figures/pnl_curve_pts_l10.png}
        \caption{Points (L10 line)}
    \end{subfigure}
    \hfill
    \begin{subfigure}[b]{0.32\textwidth}
        \includegraphics[width=\textwidth]{figures/pnl_curve_reb_l10.png}
        \caption{Rebounds (L10 line)}
    \end{subfigure}
    \hfill
    \begin{subfigure}[b]{0.32\textwidth}
        \includegraphics[width=\textwidth]{figures/pnl_curve_ast_l10.png}
        \caption{Assists (L10 line)}
    \end{subfigure}
    \caption{Cumulative flat-unit P\&L over time (2021--2024) against the
             $L_{10}$ synthetic line at a 1-unit edge threshold. Each point
             represents a single bet outcome; shading indicates profit (green)
             and loss (red) regions.}
    \label{fig:pnl}
\end{figure}

Against the $L_{10}$ line, ROI rises monotonically with threshold for all
three statistics: from $+6.3\%$ (all bets, $\tau=0$) to $+30.0\%$
($\tau=2.5$, $N=6{,}128$) for points.
This pattern is consistent with the model's edge being concentrated in
cases where its prediction diverges substantially from the naive recency
estimate --- precisely the situations where the model's richer feature set
provides the most additional information.

The EMA$_5$ line systematically yields higher ROI than the $L_{10}$ line
at the same threshold across all statistics.
This should not be interpreted as evidence that EMA$_5$ is a weaker market
proxy; rather, EMA$_5$'s shorter effective window introduces additional
noise, creating more observations where XGBoost and the EMA$_5$ line disagree
by more than $\tau$.
Many of those disagreements reflect the line being wrong (over-reacting to
a single game) rather than XGBoost being right.

The most conservative and epistemically meaningful comparison is against the
linear regression line.
Both models share identical feature sets, eliminating information asymmetry.
XGBoost achieves +13.1\% ROI on points ($N=17{,}685$), +20.8\% on rebounds
($N=1{,}525$), and +15.6\% on assists ($N=793$) at a threshold of 1.0.
Win rates of 59.3\%, 63.3\%, and 60.5\% respectively --- all well above the
52.38\% break-even --- confirm that the non-linear model consistently
identifies cases where the linear model systematically under- or
over-predicts, and that XGBoost's prediction is on the correct side of
those deviations.
The smaller sample sizes for rebounds and assists at this threshold warrant
caution in interpretation; the rebound result in particular ($N=1{,}525$)
has wider confidence intervals and should be replicated on additional data.

% =============================================================
\section{Limitations and Future Work}
\label{sec:limits}
% =============================================================

\paragraph{Synthetic market lines.}
All betting simulation results are relative to synthetically constructed
lines rather than historical sportsbook closing lines.
Real prop lines incorporate injury reports, lineup information, travel
schedules, and sharp money movement --- factors absent from our model.
The $L_{10}$, EMA$_5$, and linear regression lines are all computable
from the same dataset and do not incorporate market microstructure.
Evaluation against historical closing lines (e.g., via commercial odds
data providers such as OddsJam or The Odds API) would yield more
practically meaningful ROI estimates and is the most important
direction for future validation.

\paragraph{Injury and lineup data.}
The model does not incorporate real-time roster information.
When a star player is absent, the minutes and statistical distribution
of teammates shifts substantially.
Integrating lineup-adjusted features --- for example, estimated minutes
shares from Vegas injury reports or projected starter availability ---
is a known improvement opportunity in NBA prediction
literature~\citep{chen2021hybrid}.

\paragraph{Probability calibration.}
The probability of exceeding the line is estimated via the normal CDF
with the player's 10-game standard deviation as the dispersion parameter.
This distributional assumption is not validated against empirical
frequencies.
Calibration plots and post-hoc methods such as Platt
scaling~\citep{platt1999probabilistic} or isotonic regression are
left as future work.

\paragraph{Performance drift.}
The slight upward trend in MAE across folds (fold 1: 4.770 points,
fold 4: 4.859 points) suggests mild distribution shift over time.
Techniques such as concept drift detection, time-decayed sample
weighting, or season-specific normalisation may improve stability
in future iterations.

\paragraph{Historical era coverage.}
The training data spans 2016--2024, but league play style continues
to evolve.
Models trained on all available historical data may be distorted by
structural shifts in pace, three-point volume, and role specialisation
between early and late seasons in this window.
Restricting training to a shorter modern-era window or applying
exponential time-decay to older observations is worth exploring.

% =============================================================
\section{Conclusion}
\label{sec:conclusion}
% =============================================================

We presented a gradient boosting regression pipeline for NBA player prop
prediction that achieves consistent temporal generalization across four
held-out seasons.
Walk-forward cross-validation yields mean absolute errors of 4.84 points,
1.97 rebounds, and 1.41 assists --- improving over a naive recency baseline
by 4.5--5.0\%.
The near-identical performance of XGBoost and random forest, combined with
only modest gains over linear regression, indicates that the feature
engineering --- particularly the rolling window structure, exponential
moving average, season baseline, and opponent defensive context ---
drives the majority of predictive improvement over a simple recency heuristic.

SHAP analysis confirms a consistent three-feature hierarchy across all
statistics: the 10-game rolling mean is dominant, followed by the season
expanding mean and the EMA$_5$ momentum indicator.
Opponent defensive context contributes meaningfully as the fourth-ranked
feature for points and assists, validating its inclusion.
A notable cross-stat finding is that home/away status is uniquely important
for assists (rank 4), consistent with home-court advantages in playmaking
but not in scoring or rebounding.

A betting simulation against a linear regression synthetic line --- the most
conservative benchmark, as both models share identical features --- produces
positive ROI across all statistics (PTS: $+13.1\%$, REB: $+20.8\%$,
AST: $+15.6\%$ at a 1-unit edge threshold), demonstrating that XGBoost's
non-linear capacity captures structure a linear model does not.
Win rates of 59--63\% against the linear baseline confirm directional
accuracy that consistently exceeds the 52.38\% break-even threshold.

These results position the model as a viable first-stage filter for prop
betting analysis, with the understanding that practical deployment would
require integration with real-time injury data, historical closing lines
for honest ROI estimation, and probability calibration for reliable
expected-value calculations.

All code, trained models, and evaluation scripts are available at:
\url{https://github.com/andomo3/nbaPropsPrediction}

% =============================================================
% Bibliography
% =============================================================
\newpage
\bibliographystyle{plainnat}
\bibliography{references}

% =============================================================
\appendix
% =============================================================

\section{Per-Fold MAE Results}
\label{app:folds}

\begin{table}[H]
\centering
\caption{MAE by model, statistic, and fold year.}
\label{tab:fold_detail}
\begin{tabular}{llcccc}
\toprule
\textbf{Stat} & \textbf{Fold} & \textbf{Naive L5} & \textbf{Linear} & \textbf{RF} & \textbf{XGBoost} \\
\midrule
\multirow{4}{*}{Points}
  & 2021 & 4.999 & 4.805 & 4.767 & 4.770 \\
  & 2022 & 5.092 & 4.904 & 4.856 & 4.853 \\
  & 2023 & 5.124 & 4.928 & 4.873 & 4.877 \\
  & 2024 & 5.113 & 4.899 & 4.859 & 4.859 \\
\midrule
\multirow{4}{*}{Rebounds}
  & 2021 & 2.069 & 1.981 & 1.963 & 1.963 \\
  & 2022 & 2.085 & 2.001 & 1.985 & 1.982 \\
  & 2023 & 2.063 & 1.982 & 1.964 & 1.964 \\
  & 2024 & 2.061 & 1.979 & 1.961 & 1.957 \\
\midrule
\multirow{4}{*}{Assists}
  & 2021 & 1.431 & 1.381 & 1.368 & 1.366 \\
  & 2022 & 1.467 & 1.419 & 1.403 & 1.403 \\
  & 2023 & 1.488 & 1.439 & 1.424 & 1.423 \\
  & 2024 & 1.528 & 1.475 & 1.458 & 1.458 \\
\bottomrule
\end{tabular}
\end{table}

\section{SHAP Feature Importance Tables}
\label{app:shap}

\begin{table}[H]
\centering
\caption{Mean absolute SHAP values --- Points model (top 10 features,
         2024 season hold-out, $N=5{,}000$ samples).}
\label{tab:shap_pts}
\begin{tabular}{lc}
\toprule
\textbf{Feature} & \textbf{Mean $|\text{SHAP}|$} \\
\midrule
pts\_L10              & 2.752 \\
season\_avg\_pts      & 1.134 \\
pts\_ema\_L5          & 1.100 \\
opp\_pts\_allowed\_L10 & 0.346 \\
min\_L5               & 0.313 \\
fg\_pct\_L10          & 0.221 \\
fg\_pct\_L5           & 0.212 \\
is\_home              & 0.187 \\
days\_rest            & 0.145 \\
min\_L10              & 0.134 \\
\bottomrule
\end{tabular}
\end{table}

\begin{table}[H]
\centering
\caption{Mean absolute SHAP values --- Rebounds model (top 10 features,
         2024 season hold-out, $N=5{,}000$ samples).}
\label{tab:shap_reb}
\begin{tabular}{lc}
\toprule
\textbf{Feature} & \textbf{Mean $|\text{SHAP}|$} \\
\midrule
reb\_L10              & 0.931 \\
season\_avg\_reb      & 0.643 \\
reb\_ema\_L5          & 0.367 \\
min\_L5               & 0.185 \\
min\_L10              & 0.159 \\
is\_home              & 0.105 \\
opp\_reb\_allowed\_L10 & 0.075 \\
days\_rest            & 0.063 \\
reb\_L5               & 0.051 \\
pts\_L5               & 0.038 \\
\bottomrule
\end{tabular}
\end{table}

\begin{table}[H]
\centering
\caption{Mean absolute SHAP values --- Assists model (top 10 features,
         2024 season hold-out, $N=5{,}000$ samples).}
\label{tab:shap_ast}
\begin{tabular}{lc}
\toprule
\textbf{Feature} & \textbf{Mean $|\text{SHAP}|$} \\
\midrule
ast\_L10              & 0.770 \\
season\_avg\_ast      & 0.406 \\
ast\_ema\_L5          & 0.247 \\
is\_home              & 0.093 \\
opp\_ast\_allowed\_L10 & 0.089 \\
min\_L5               & 0.043 \\
pts\_L5               & 0.032 \\
ast\_L5               & 0.025 \\
days\_rest            & 0.021 \\
hot\_cold\_ast        & 0.009 \\
\bottomrule
\end{tabular}
\end{table}

\section{Full Betting Simulation Results}
\label{app:betting}

\begin{table}[H]
\centering
\caption{Flat-unit betting simulation: ROI and win rate by statistic,
         synthetic line type, and edge threshold $\tau$.
         Break-even win rate at $-110$ juice $= 52.38\%$.}
\label{tab:betting_full}
\small
\begin{tabular}{llccccc}
\toprule
\textbf{Stat} & \textbf{Line} & \textbf{$\tau$} & \textbf{N bets} & \textbf{N won} & \textbf{Win rate} & \textbf{ROI} \\
\midrule
\multirow{6}{*}{Points}
  & \multirow{6}{*}{$L_{10}$}
    & 0.0 & 94,462 & 52,603 & 55.7\% & $+0.063$ \\
  & & 0.5 & 67,242 & 38,690 & 57.5\% & $+0.099$ \\
  & & 1.0 & 42,918 & 25,689 & 59.9\% & $+0.143$ \\
  & & 1.5 & 24,677 & 15,390 & 62.4\% & $+0.191$ \\
  & & 2.0 & 12,629 &  8,259 & 65.4\% & $+0.249$ \\
  & & 2.5 &  6,128 &  4,172 & 68.1\% & $+0.300$ \\
\midrule
\multirow{6}{*}{Points}
  & \multirow{6}{*}{EMA$_5$}
    & 0.0 & 94,462 & 55,184 & 58.4\% & $+0.115$ \\
  & & 0.5 & 74,117 & 44,755 & 60.4\% & $+0.153$ \\
  & & 1.0 & 55,410 & 34,761 & 62.7\% & $+0.198$ \\
  & & 1.5 & 39,394 & 25,628 & 65.1\% & $+0.242$ \\
  & & 2.0 & 26,535 & 17,876 & 67.4\% & $+0.286$ \\
  & & 2.5 & 17,087 & 11,889 & 69.6\% & $+0.328$ \\
\midrule
\multirow{6}{*}{Points}
  & \multirow{6}{*}{LinReg}
    & 0.0 & 94,462 & 51,077 & 54.1\% & $+0.032$ \\
  & & 0.5 & 46,364 & 26,112 & 56.3\% & $+0.075$ \\
  & & 1.0 & 17,685 & 10,480 & 59.3\% & $+0.131$ \\
  & & 1.5 &  6,219 &  3,797 & 61.1\% & $+0.166$ \\
  & & 2.0 &  2,308 &  1,454 & 63.0\% & $+0.203$ \\
  & & 2.5 &    961 &    624 & 64.9\% & $+0.240$ \\
\midrule
\multirow{6}{*}{Rebounds}
  & \multirow{6}{*}{$L_{10}$}
    & 0.0 & 94,462 & 51,920 & 55.0\% & $+0.049$ \\
  & & 0.5 & 29,471 & 17,738 & 60.2\% & $+0.149$ \\
  & & 1.0 &  6,416 &  4,201 & 65.5\% & $+0.250$ \\
  & & 1.5 &  1,488 &  1,030 & 69.2\% & $+0.322$ \\
  & & 2.0 &    430 &    319 & 74.2\% & $+0.416$ \\
  & & 2.5 &    125 &     92 & 73.6\% & $+0.405$ \\
\midrule
\multirow{6}{*}{Rebounds}
  & \multirow{6}{*}{LinReg}
    & 0.0 & 94,462 & 51,464 & 54.5\% & $+0.040$ \\
  & & 0.5 & 13,083 &  7,772 & 59.4\% & $+0.134$ \\
  & & 1.0 &  1,525 &    965 & 63.3\% & $+0.208$ \\
  & & 1.5 &    262 &    180 & 68.7\% & $+0.312$ \\
\midrule
\multirow{6}{*}{Assists}
  & \multirow{6}{*}{$L_{10}$}
    & 0.0 & 94,462 & 50,695 & 53.7\% & $+0.025$ \\
  & & 0.5 & 16,397 & 10,057 & 61.3\% & $+0.171$ \\
  & & 1.0 &  2,552 &  1,731 & 67.8\% & $+0.295$ \\
  & & 1.5 &    500 &    357 & 71.4\% & $+0.363$ \\
\midrule
\multirow{4}{*}{Assists}
  & \multirow{4}{*}{LinReg}
    & 0.0 & 94,462 & 51,998 & 55.1\% & $+0.051$ \\
  & & 0.5 &  7,432 &  4,325 & 58.2\% & $+0.111$ \\
  & & 1.0 &    793 &    480 & 60.5\% & $+0.156$ \\
  & & 1.5 &     99 &     69 & 69.7\% & $+0.331$ \\
\bottomrule
\end{tabular}
\end{table}

\end{document}
